\documentclass[10pt,twocolumn,letterpaper]{article}

\usepackage[review]{cvpr}
\usepackage{times}
\usepackage{epsfig}
\usepackage{graphicx}
\usepackage{amsmath}
\usepackage{amssymb}
\usepackage{booktabs}
\usepackage{tabularx}
\usepackage[pagebackref,breaklinks,colorlinks]{hyperref}
\usepackage[capitalize]{cleveref}

\def\cvprPaperID{0000}
\def\confName{CVPR}
\def\confYear{2026}

\begin{document}

\title{Lightweight Temporal Modeling for Driver Visual Saliency Prediction}

\author{
Author Name\\
Department and Institution\\
{\tt\small author@example.com}
}
\maketitle

\begin{abstract}
This work studies current-frame visual-saliency prediction for panoramic
driving videos using the DriverGaze360 dataset. Given a short sequence of RGB
frames ending at time \(t\), the task is to predict the dataset-provided
saliency map associated with frame \(t\). We implement a lightweight model
based on a MobileNetV3-Large encoder, Temporal Shift Modules (TSM), temporal
feature aggregation, and a multi-scale decoder. A single-frame model is
defined as a baseline for evaluating whether temporal information contributes
to prediction quality. The project also documents frame alignment checks and
aspect-preserving preprocessing for the panoramic images. At the time of
writing, a complete training and evaluation pipeline has not been validated;
therefore, we report no quantitative results and make no claim of
improvement. The paper specifies the task, implementation, and controlled
experiments required to test the temporal-modeling hypothesis.
\end{abstract}

\section{Project Team and Completion}
\label{sec:team}

The member names, task assignments, and completion estimates were not
provided in the project materials available for this draft. The table below
is a reporting template and must be completed by the team; no individual
contribution or completion percentage is inferred here.

\begin{table}[t]
    \centering
    \caption{Team members and self-reported task completion.}
    \label{tab:team}
    \small
    \begin{tabularx}{\linewidth}{l X c}
        \toprule
        Member & Assigned work / evidence & Completion \\
        \midrule
        To be completed & To be completed & -- \\
        To be completed & To be completed & -- \\
        To be completed & To be completed & -- \\
        \bottomrule
    \end{tabularx}
\end{table}

\section{Introduction}
\label{sec:intro}

Visual attention is relevant to understanding driver behavior and to
developing driver-assistance systems. A visual-saliency map represents a
spatial distribution of attention over an image and offers one way to analyze
which regions of a driving scene attract gaze. In panoramic driving scenes,
the model must process a wide field of view while preserving spatial detail.
Studying this task is useful for evaluating whether models can represent
driver-attention patterns across a panoramic scene. A saliency prediction
alone, however, does not establish why a driver looked at a region or whether
the behavior was safe.

This work considers a focused prediction task on DriverGaze360
\cite{govil2026drivergaze360}. Given a short RGB clip ending at frame \(t\),
the model predicts the dataset-provided saliency map for that same frame.
This is current-frame prediction, not future gaze forecasting. Our research
question is:
\begin{quote}
Does lightweight temporal feature modeling improve current-frame saliency
prediction compared with a single-frame image model?
\end{quote}

We implement a MobileNetV3-Large encoder \cite{howard2019mobilenetv3} and use
TSM \cite{lin2019tsm} to exchange a fraction of feature channels across
neighboring frames. A multi-scale decoder produces a dense saliency map.
These design choices motivate a controlled comparison, but do not by
themselves establish that temporal modeling improves accuracy. We therefore
distinguish the implemented architecture from the experiments that remain to
be run.

\paragraph{Contributions.}
This project makes the following implementation and evaluation contributions:
\begin{itemize}
    \item It defines a current-frame saliency-prediction task using panoramic
    RGB clips and the corresponding final-frame saliency annotation.
    \item It implements a lightweight MobileNetV3-based model with optional
    temporal feature exchange and multi-scale decoding, together with a
    single-frame baseline.
    \item It specifies data-alignment checks and controlled ablations needed
    to evaluate the contribution of temporal modeling without presuming a
    performance gain.
\end{itemize}

\section{Method}
\label{sec:method}

\subsection{Problem statement, input, and output}

Let \(I_t\) denote the RGB frame at time \(t\), and let \(S_t\) denote its
provided saliency map. Given a clip of \(T\) frames ending at \(t\), the model
predicts the saliency map for the final input frame:
\begin{equation}
    X_t = [I_{t-(T-1)}, \ldots, I_{t-1}, I_t],
    \qquad
    \hat{S}_t = f_\theta(X_t).
    \label{eq:task}
\end{equation}
This is current-frame prediction, not future gaze forecasting. The target is
the dataset-provided saliency annotation paired with \(I_t\).

\paragraph{Input.}
For the default temporal model, the input is \(T=16\) RGB frames with
preprocessed tensor shape \([B,16,3,224,1120]\), where \(B\) is the batch
size. The single-frame baseline uses the final frame \(I_t\) from this clip.

\paragraph{Output.}
The model predicts one saliency map for the final frame, with shape
\([B,1,224,1120]\). The temporal model returns logits by default; the
single-frame baseline applies sigmoid and returns values in \([0,1]\).
Predictions must be converted to the representation expected by the selected
loss and evaluation metrics.

\paragraph{Illustrative example.}
\begin{figure}[t]
    \centering
    \fbox{\parbox{0.92\linewidth}{
        \centering
        RGB clip \(\{I_{t-15},\ldots,I_t\}\)
        \(\longrightarrow\)
        saliency model
        \(\longrightarrow\)
        predicted map \(\hat{S}_t\)\\[2pt]
        \small Schematic only; replace with a real input, ground-truth map,
        and model prediction after inference has been run.
    }}
    \caption{Conceptual input-output example for current-frame saliency
    prediction. This schematic contains no experimental prediction.}
    \label{fig:io-example}
\end{figure}

\subsection{Architecture}

The implemented \texttt{DriverGazeNet} contains a MobileNetV3-Large image
encoder, optional TSM modules, per-stage temporal pooling, and a lightweight
multi-scale saliency decoder. The encoder exposes features from four stages
with channel dimensions \((24, 40, 112, 960)\). The decoder projects these
features to a shared width, combines them in a top-down sequence, refines the
fused features with depthwise-separable convolutions, and upsamples to the
input resolution.

\subsection{Temporal Shift Module}

TSM exchanges feature channels across neighboring frames without introducing
a conventional temporal convolution \cite{lin2019tsm}. In this
implementation, the channel tensor is divided into groups: one group is
shifted from the next frame, a second group from the previous frame, and the
remaining channels stay at the current frame. The operation is inserted into
residual blocks in the MobileNet feature extractor. At temporal boundaries,
the shifted groups are zero-filled.

\subsection{Temporal aggregation and output}

After feature extraction, temporal pooling is applied independently to each
selected encoder stage. The implementation supports last-frame selection,
mean pooling, and learned attention pooling. The default configuration uses
last-frame pooling. The decoder returns one saliency logit map at the input
resolution; sigmoid may be applied for probability-map output.

The single-frame baseline uses only \(I_t\), disables TSM, and applies sigmoid
in the decoder. The temporal model and baseline should be trained with losses
that match their output conventions: a probability output must not receive a
second sigmoid operation inside the loss.

\section{Experimental Evaluation}
\label{sec:experiments}

\subsection{Dataset}

DriverGaze360 provides panoramic driving recordings with gaze-related
annotations and additional modalities, including saliency, depth, and
instance-segmentation data \cite{govil2026drivergaze360}. The original release
stores videos and metadata. The official preparation script extracts video
frames as numbered image files and unpacks instance-segmentation archives.
Frame filenames are one-indexed to align with the \texttt{Frame} field in
\texttt{sim\_gaze\_df.csv}.

The official project supports broader object-level attention tasks. In
contrast, this work uses RGB clips and the associated saliency map only. It
does not train on depth or instance-segmentation targets and does not
reconstruct a new saliency target from gaze coordinates.

\subsection{Sampling and frame alignment}

For the default settings, \(T=16\) and the frame stride is one. DriverGaze360
videos are sampled at 30 frames per second, so the clip spans 0.5 seconds
between its first and final frames. The default clip-start step is eight
frames (approximately 0.267 seconds), and the loader skips the first 30
frames of each recording.

The official preparation script uses FFmpeg to extract RGB, saliency, and
depth videos into numbered JPEG frames. This representation allows the
loader to read selected frames by index and construct overlapping clips.
The project-specific \texttt{DG360LiteDataset} checks that CSV frame IDs are
integer-valued and strictly increasing, and that corresponding RGB and
saliency filenames exist. It uses those frame IDs when loading samples.

For the locally audited subset of 46 recordings, CSV frame IDs matched the
RGB and saliency filenames, and inspected source images had dimensions of
\(6400 \times 720\). These checks establish file-level alignment and observed
shape consistency for the audited subset; they do not establish the semantic
validity of every saliency annotation.

\subsection{Panorama resizing}

The loader resizes RGB and saliency images using the same
aspect-preserving fit-and-center-pad transform to a \(224 \times 1120\)
\((H \times W)\) canvas. For a \(6400 \times 720\) source image, the resized
content is approximately \(1120 \times 126\) pixels \((W \times H)\), with
zero padding above and below. RGB values are scaled to floating point and
ImageNet-normalized; saliency values are scaled to \([0,1]\), but are not
normalized to sum to one by the loader.

The padding creates pixels outside the source image. Losses and metrics
should therefore either exclude those pixels with a valid-image mask or
explicitly justify treating them as background. This choice has not yet been
validated in a complete training and evaluation pipeline.

\subsection{Experimental setup}

\subsubsection{Data split}

The first experiments should preserve the dataset's original training and
validation partitions. The project download script selects approximately
20\% of recordings independently within each source split using a fixed
random seed. This recording-level subset does not guarantee that drivers are
disjoint across partitions. Accordingly, results from this setup must not be
described as generalization to unseen drivers. A participant-grouped split
would be required to support that claim.

\subsubsection{Baselines and ablations}

All models should use identical partitions, preprocessing, target
normalization, and evaluation procedures. The planned comparisons are:
\begin{enumerate}
    \item A spatial-prior baseline, such as the mean training-set saliency map.
    \item A single-frame MobileNetV3 model using only \(I_t\).
    \item A 16-frame model with temporal aggregation but without TSM.
    \item A 16-frame model with TSM.
\end{enumerate}

Temporal pooling modes (last, mean, or learned attention) should be treated as
separate controlled ablations if they are evaluated. Selection of model
variants and hyperparameters must use validation data only.

\subsubsection{Training setup and evaluation metrics}

The training report should state the optimizer, learning rate, batch size,
epoch count, random seed, hardware, loss definition, and target
normalization used in the actual run. These settings have not yet been
confirmed by a completed experiment and are therefore not specified here as
measured protocol.

We plan to evaluate saliency predictions using Kullback--Leibler divergence
(KLD), correlation coefficient (CC), similarity (SIM), and normalized
scanpath saliency (NSS). Metric definitions and prediction/target
normalization must be reported explicitly. NSS additionally requires a
fixation map that is appropriate for the evaluated frame. If efficiency
claims are made, model size and inference latency should also be measured
and reported with hardware details.

\subsubsection{Experimental environment}

The hardware and software environment for a completed training or inference
run has not been recorded in the verified project state. Before reporting
results, record the GPU/CPU model, memory, operating system, Python and
PyTorch versions, CUDA version (if applicable), input resolution, batch size,
and inference timing procedure. These values are intentionally left
unspecified rather than inferred from the development environment.

\subsection{Results and current status}
\label{sec:results}

No complete training run or quantitative evaluation has been validated in
the current project state. Consequently, no KLD, CC, SIM, NSS, latency, or
model-size results are claimed. The verified work currently consists of the
project data loader, frame-pairing checks on a local subset, preview
functionality, and model forward-pass implementations. Training, checkpoint
selection, and metric behavior still require end-to-end validation.

\subsubsection{Quantitative results}

The following table is a results template, not a report of measured
performance. Fill it only after running the stated experiments with the
same split, preprocessing, and metric definitions.

\begin{table}[t]
    \centering
    \caption{Planned experiments. Quantitative results have not yet been
    measured.}
    \label{tab:planned}
    \begin{tabular}{lccccc}
        \toprule
        Experiment & KLD & CC & SIM & NSS & Status \\
        \midrule
        Spatial prior & -- & -- & -- & -- & Not run \\
        Single-frame & -- & -- & -- & -- & Not run \\
        16-frame, no TSM & -- & -- & -- & -- & Not run \\
        16-frame, with TSM & -- & -- & -- & -- & Not run \\
        \bottomrule
    \end{tabular}
\end{table}

\subsubsection{Qualitative results}

No verified model predictions are available yet. After inference, replace the
boxes below with a representative RGB frame, its dataset-provided saliency
target, and the corresponding prediction. Select examples from the evaluation
split, and do not choose examples solely because they make the model appear
successful.

\begin{figure}[t]
    \centering
    \setlength{\fboxsep}{5pt}
    \fbox{\parbox[c][2.6cm][c]{0.27\linewidth}{\centering RGB input\\[2pt]
    \small Replace with real frame}}
    \hfill
    \fbox{\parbox[c][2.6cm][c]{0.27\linewidth}{\centering Ground truth\\[2pt]
    \small Replace with paired saliency}}
    \hfill
    \fbox{\parbox[c][2.6cm][c]{0.27\linewidth}{\centering Prediction\\[2pt]
    \small Replace after inference}}
    \caption{Qualitative comparison template. The boxes are placeholders and
    are not model outputs.}
    \label{fig:qualitative-template}
\end{figure}

\section{Limitations and Future Work}
\label{sec:limitations}

\subsection{Strengths of the current design}

The project has a clearly scoped current-frame prediction target, a
single-frame baseline for comparison, and a lightweight temporal mechanism
whose contribution can be tested by ablation. The loader also checks
frame-ID alignment and preserves the panoramic aspect ratio during resizing.
These are design and implementation properties, not evidence of improved
predictive performance.

\subsection{Limitations}

This paper describes an implementation and experimental plan rather than a
completed empirical study. The local subset is only a fraction of the full
dataset, and recording-level sampling does not ensure participant-level
separation. File alignment and shape checks are necessary data-quality steps,
but cannot verify that each target map semantically represents driver gaze.

The panorama resize-and-pad operation also introduces regions outside the
original image. The training loss and metrics should handle these pixels
consistently, ideally with an explicit valid-image mask. The immediate next
steps are to validate one forward/backward pass, overfit a small batch as a
training sanity check, implement and validate the loss and metrics, and run
the spatial-prior and single-frame baselines before the temporal ablations.
Any claims about generalization to unseen drivers require a separate
participant-held-out evaluation.

\section{Conclusion}
\label{sec:conclusion}

We define a current-frame driver visual-saliency prediction task on panoramic
video and describe a lightweight MobileNetV3-based model with optional TSM
and multi-scale decoding. A single-frame baseline provides the reference
needed to test whether temporal information is useful. Since the training
and evaluation pipeline has not yet been validated, the effectiveness of the
temporal components remains an open empirical question.

\bibliographystyle{ieeenat_fullname}
\bibliography{references}

\end{document}
